\section{Fresh Workflow Generalization and Gate Attribution}
\label{app:generalization-gate}

This separately frozen follow-up extends the adapted AWM comparison in
\cref{app:matched-reuse} to three workflow families and two model strings:
\code{gpt-5.6-luna} and \code{gpt-6-luna}. All records come from the same pinned
repository snapshot. Thus the study tests workflow and model transfer, not
cross-domain or cross-provider generalization.

\paragraph{Prospective selection and execution.}
All record identifiers in prior retained studies are excluded. Backlog tasks
are allocated first, followed by issue triage and pull-request outcome tasks,
with fixed hash ranking within source-only eligibility. Titles must be
representable by the answer schema, and issue tasks must expose a source
excerpt meeting the minimum length. No successful dispatch or provider answer
is an eligibility criterion. Six pilot records per family precede 144 backlog,
600 issue, and 600 pull-request primary records. Families are disjoint; models
share records. The design schedules 8,172 episodes, including 108 pilot
episodes, with all six three-arm permutations balanced within each cell.
Models alternate order by record. Eight record groups may execute concurrently;
all arms and models within a record execute sequentially.

The primary class mix is 169 bug, 59 enhancement and 372 other issues;
517 merged, 82 closed-unmerged and one open pull request; and 109 discussed
unowned plus 35 awaiting-response backlog records. There are no owned backlog
records in the fresh cohort. The hash-selected residual population therefore
does not supply balanced branch coverage, and the open/owned branches cannot
support generalization claims from this follow-up. The source audit finds no
structural title/excerpt/owner representability defect, but flags 27 issue
records with nonliteral bug labels; exact-label errors can differ from semantic
classification errors. This is an author audit, not independent adjudication.

The baseline, GAC, and offline adapted AWM arms share each family's prompt,
source tools, output schema and factual grader. GAC uses the existing SDK model
wrapper for issue triage and the outer pre-model composite runner for backlog
and pull-request tasks; refusals use the unchanged agent. These placements are
part of the treatment. Historical discovery is reconstructed under each model's
manifest, without model-specific rediscovery. AWM induces one memory per cell
from that family's eligible historical action traces, using the pinned adapter
in \cref{app:matched-reuse}; the memory is frozen before pilots, without manual
rewriting or online updates. The reused adapter retains its generic
``Software,Issue triage,github'' induction header for every family; the
action examples and tool schemas are family-specific, and every retained
memory names that family's three tools. This transport/domain adaptation is
not native AWM evaluation. GAC retains its historical split. Induction and
reconstruction costs are reported separately from deployment.

Settings are low reasoning/verbosity, no parallel tools, 2,048 output tokens,
eight SDK turns, a 120-second timeout and zero SDK/HTTP retries. Pilot continuation
requires execution, schema and usage completeness and actual compaction in
every cell; favorable accuracy is not a continuation criterion. Failed and
unattempted scheduled tasks remain incorrect in accuracy denominators. There
are no replacement tasks or repeat attempts.

\paragraph{Accuracy margin and precision.}
The primary contrast is GAC minus adapted AWM on exact source-grounded answer
fields; tool-order compliance is reported separately. The registered
noninferiority margin is five percentage points, a research tolerance rather
than an operational safety threshold. For each cell, let $L$ count
AWM-correct/GAC-incorrect pairs among $n$ scheduled records. A one-sided
Clopper--Pearson upper bound $U$ on $L/n$ at error probability $0.05/6$ yields
a conservative lower bound $-U$ on the accuracy difference, since it ignores
GAC-only successes. Noninferiority requires all scheduled pairs attempted and
$U<0.05$. Bonferroni controls the six comparisons under the record-sampling
assumptions without requiring model independence. Models are never counted as
additional independent records in a pooled interval.

At $n=600$, at most 17 losses pass; power is 93.9\% under a 2\% loss
probability, but 46.7\% under a 3\% probability. At $n=144$, only zero or one
loss passes, and corresponding power is 21.5\% or 6.8\%. The backlog cell is
therefore a low-precision replication, not an adequately powered test under
those alternatives. Tighter one- and two-point margins are sensitivity analyses.
Resource intervals use 10,000 paired bootstrap draws on completed pairs;
all-attempted resource totals also retain observed failure usage. Concurrent
wall-clock measurements include shared runner/checkpoint overhead. Missing usage
is never interpreted as zero cost. Task/author/time independence is unproven.

\paragraph{Gate attribution.}
Before paid execution, the real dispatcher runs each selected record under
full statistical admission, support-only admission, and deterministic guards.
Support-only requires the retained 92-group calibration floor but disables the
score and risk admission; deterministic guards also remove the support-count
check. Both ablations preserve the exact candidate program, hard guards,
verifier, effect checks and continuation contract. They are uncertified research
ablations on already admitted candidates, not alternative certified compilers.

All three paths produce identical decisions, calls and continuation evidence on
all 2,724 record/model cases, including pilots. Recorded gate scores are constant
within each family on both models: 0.1237 for backlog, 0.0815 for issue triage,
and 0.1308 for pull requests. The selected entry features therefore provide no
observed risk ranking within these cohorts. Backlog and pull-request
candidates dispatch throughout; issue candidates dispatch on 600 of 606 records
per model and refuse the same six records under every variant. Consequently,
only one live GAC path is run: duplicate stochastic answers would not constitute
independent evidence that identical routing policies differ. This is a null
result for incremental runtime gate value on these candidates. It neither
establishes universal gate equivalence nor measures whether statistical admission
excluded bad candidates during compilation. The earlier learned end-to-end gate
retirement remains unchanged.

\paragraph{Completed results.}
All 8,172 scheduled episodes were attempted: 8,171 completed and one baseline
PR episode on \code{gpt-5.6-luna} failed with a provider server error. It remains
incorrect, with no retry. All 2,724 live GAC decisions match their frozen
provider-free probes, with zero dispatch mismatches. The pilot has
107/108 correct answers; primary results are separate in
\cref{tab:generalization-accuracy,tab:generalization-resources}.

\begin{table}[htbp]
\centering\small
\begin{tabular}{llrrrrrr}
\toprule
Family & Model & $n$ & GAC & AWM & Base & $100U$ & NI \\
\midrule
Backlog & 5.6 Luna & 144 & 143 & 142 & 144 & 3.27 & yes \\
Backlog & 6 Luna & 144 & 142 & 143 & 144 & 5.86 & no \\
Issue triage & 5.6 Luna & 600 & 584 & 583 & 590 & 3.61 & yes \\
Issue triage & 6 Luna & 600 & 583 & 583 & 582 & 4.04 & yes \\
PR outcome & 5.6 Luna & 600 & 597 & 598 & 597 & 1.43 & yes \\
PR outcome & 6 Luna & 600 & 598 & 600 & 599 & 1.43 & yes \\
\bottomrule
\end{tabular}

\caption{Primary exact-answer successes over scheduled records. AWM denotes
the offline API-task adaptation. $100U$ is the simultaneous conservative
accuracy-loss bound in percentage points, not an observed accuracy difference.
NI uses the registered five-point margin.}
\label{tab:generalization-accuracy}
\end{table}

Five cells pass the five-point margin. Second-model backlog has two AWM-only
successes and one GAC-only success; its 5.86-point bound does not establish
noninferiority or inferiority. Only the two PR cells pass a two-point margin,
and none passes one point. Thus this study does not establish exact accuracy
equality or a uniformly small quality loss across all cells.

\begin{table}[htbp]
\centering\small
\setlength{\tabcolsep}{3pt}
\begin{tabular}{llrrrrr}
\toprule
Family & Model & GAC req. & AWM req. & GAC USD/1k & AWM USD/1k & Saving \\
\midrule
Backlog & 5.6 Luna & 1.00 & 4.00 & 0.204 & 0.806 & 74.7\% \\
Backlog & 6 Luna & 1.00 & 4.00 & 0.088 & 0.380 & 76.8\% \\
Issue triage & 5.6 Luna & 2.00 & 4.00 & 0.588 & 0.931 & 36.9\% \\
Issue triage & 6 Luna & 2.00 & 4.00 & 0.210 & 0.411 & 48.9\% \\
PR outcome & 5.6 Luna & 1.00 & 4.00 & 0.164 & 0.739 & 77.8\% \\
PR outcome & 6 Luna & 1.00 & 4.00 & 0.074 & 0.353 & 79.1\% \\
\bottomrule
\end{tabular}

\caption{Observed deployment resources per attempted primary episode;
USD/1k scales estimated cost to 1,000 episodes, excluding induction. Savings
are relative to adapted AWM. Both compared arms have complete usage in every
cell. Frozen API prices are estimates, not provider invoices.}
\label{tab:generalization-resources}
\end{table}

Paired cost and latency differences favor GAC in every cell; the descriptive
95\% intervals exclude zero, with 144 or 600 complete pairs. Tool-call means
are equal between the compared arms. The full resource intervals are retained
in \path{analysis.json}; latency includes shared runner overhead. New observed
cost, including six inductions, is \$3.864269; induction alone is \$0.011211.
One baseline request has no returned usage; its unknown charge is not treated
as zero. The aggregate estimate plus prior and new failure reserves is
\$13.441321 against the authorized \$400 cap. No baseline-cost superiority
claim depends on the missing charge.

The post-registration author audit retains all 113 incorrect or failed
episodes, including one pilot error. Overlapping failed checks include 50
issue-excerpt, 45 category, 46 evidence-label, 12 comment-grounding and five
title checks; 45 error episodes involve nonliteral bug labels. No labels or
outcomes are changed. Author/day cluster bootstraps are descriptive sensitivity,
not replacement primary tests: the issue cohort has 386 authors (largest
cluster 70 records), and the PR cohort 190 (largest 120). Their intervals do
not establish independent sampling or repair the inconclusive backlog cell.
Per-class outcomes and every failed answer are retained in
\path{error_audit.json}. No independent human adjudication is claimed.

\paragraph{Artifacts and boundaries.}
The prospective protocol, code/environment/source hashes, frozen candidates,
cohorts, schedule, per-record gate decisions, induction memories, checkpoints
and analysis are retained under \path{experiments/generalization_gate/}.
Model aliases may drift; resolved immutable provider versions are not retained.
This evidence remains separate from the historical 810-episode six-study
summary. The study does not establish native AWM benchmark reproduction,
independent human adjudication, total lifetime optimization cost, a new risk
certificate, or cross-domain superiority.
